% META: {"id":"1SPE-PROBCOND-CO-033","chapitre":"1SPE-PROBA-COND","type_objet":"corrige","exercice_ref":"1SPE-PROBCOND-EX-033","capacites_codes":["C4"],"sources_inspiration":[],"status":"approved","origin":"generated","status_history":["generated","audited","accepted","approved"],"release_acceptance":"RELEASE_OWNER_BATCH_ACCEPTANCE","acceptance_closure_digest":"sha256:9b3ccf9a81c5520fbb7e03b7b2d3e7bf2057a02834b6d908bf3be5f49f3b553f","acceptance_receipt_digest":"sha256:4fad86f0282e0fa4e0419cca1ad6379ae7af5a280c04a4a90880f90a1c044242"}
% BEGIN-VERIFY
% from sympy import Rational
% assert Rational(1,2)*Rational(1,6) == Rational(1, 12)
% END-VERIFY
\begin{corrige}{1SPE-PROBCOND-EX-033}

\textbf{1.} $P(A) \times P(B) = \frac{1}{2} \times \frac{1}{6} = \frac{1}{12} = P(A \cap B)$. Donc $A$ et $B$ sont indépendants.

\textbf{2.} Si $A$ et $B$ sont indépendants, alors $\overline{A}$ et $B$ le sont aussi (propriété du cours). Vérification : $P(\overline{A} \cap B) = P(B) - P(A \cap B) = \frac{1}{6}-\frac{1}{12} = \frac{1}{12}$ et $P(\overline{A}) \times P(B) = \frac{1}{2} \times \frac{1}{6} = \frac{1}{12}$. \checkmark

\end{corrige}
